\scanpage{56}
would show that the conjugacy class of the sections determines the conjugacy class of $L_i$, hence $i$.

\uline{Conjecture A} Let $U$ be an algebraic anabelian geometrically connected curve over a field $K$ finitely generated over $\mathbb{Q}$. Then every section of (1) is \deleted{\unsure{defined|defined}} of one of the two preceding types, i.e. either defined by a point of $U(K)$, or by a section of an extension (2), with $i\in I(\bar K)^\Gamma$ i.e $i$ a point of $\widehat U\setminus U$, rational over $K$.
\margin{There are \inserted{\unsure{in addition|many} \unsure{some|of}} more \illegible{1} the refinements $\pi^{\Gamma^\circ}=(1)$ (first case) $\pi^{\Gamma^\circ}=L_i$ (second case)}

If one admits this conjecture, then one would conclude from it, by passage to the limit, considering the function field $L$ of $U$ and $E_L\longrightarrow \unsure{E|B}_K$ ($E_L$ can \unsure{be|be} considered as an ``infinite'' group with loops since an infinity of \inserted{classes of} loop subgroups\inserted{ $L_i$}\ldots) that every\deleted{\unsure{e}} splitting of this extension comes from a splitting of an extension of type (2), with $i\in I^\Gamma=\unsure{X}(K)$ ($X=\widehat U$). The \inserted{\unsure{conjugacy}} classes of such splittings
\margin{55}
